\documentclass[12pt,a4paper]{article}

\usepackage[utf8]{inputenc}
\usepackage{amsmath}
\usepackage{geometry}
\usepackage{setspace}
\usepackage[
    colorlinks=true,
    linkcolor=black,
    urlcolor=blue
]{hyperref}

\geometry{a4paper, margin=2cm}
\onehalfspacing

\begin{document}
\pagenumbering{gobble}

\begin{titlepage}
    \centering

    {\large Universidade Federal de São João del-Rei -- UFSJ}\\
    {\large Departamento de Ciência da Computação -- DCOMP}\\
    {\large Computação Paralela}\\

    \vspace*{5cm}

    {\Large
        \textbf{\MakeUppercase{
            Paralelização do Jogo da Vida de Conway com Pthreads
        }}
    }\\[2cm]

    {\normalsize
        Autor:\\[0.3cm]
        Samuel de Souza Carvalho Almeida
    }

    \vfill

    {\normalsize Professor: Rafael Sachetto Oliveira}\\[1cm]
    {\large São João del-Rei -- MG}\\
    {\large 2026}
\end{titlepage}

\pagenumbering{roman}
\tableofcontents
\newpage
\pagenumbering{arabic}

\section{Introdução}

Este trabalho apresenta o desenvolvimento de versões sequencial e paralela do
Jogo da Vida de Conway. A versão paralela utiliza Pthreads e o modelo de
memória compartilhada para distribuir a atualização da matriz entre diferentes
threads.

O desenvolvimento compreende o perfil da versão sequencial, a identificação
das oportunidades de paralelização, a validação funcional entre as duas versões
e uma avaliação de desempenho com diferentes tamanhos de matriz, quantidades de
gerações e números de participantes. A avaliação foi planejada para separar o
custo introduzido pela infraestrutura Pthreads do ganho obtido com o aumento do
paralelismo.

Nesta etapa, o relatório descreve a implementação e o protocolo experimental.
Os resultados quantitativos e sua discussão serão incorporados somente após a
conclusão das campanhas de experimentos.

\section{Objetivos}

O objetivo geral é desenvolver e avaliar uma versão paralela do Jogo da Vida
com Pthreads, preservando o comportamento da implementação sequencial e
utilizando os recursos de processamento disponíveis de forma controlada.

Os objetivos específicos são:

\begin{itemize}
    \item identificar a região computacionalmente relevante da versão
    sequencial;
    \item distribuir o domínio sem sobreposição de escritas e sincronizar as
    gerações corretamente;
    \item fazer a thread principal participar do cálculo, evitando que permaneça
    ociosa enquanto aguarda as threads auxiliares;
    \item verificar a equivalência das saídas sequencial e paralela; e
    \item medir speedup, eficiência, custo e sobrecusto em ambientes com
    diferentes quantidades de processadores disponíveis.
\end{itemize}

A análise posterior buscará responder em quais cargas a versão Pthreads supera
a referência sequencial, como o desempenho evolui até e além da quantidade de núcleos físicos disponíveis
e quanto do tempo paralelo decorre da
infraestrutura de criação e sincronização das threads.

\section{Jogo da Vida de Conway}

O Jogo da Vida é um autômato celular representado por uma matriz bidimensional.
Cada célula pode estar viva ou morta, e seu estado na próxima geração depende
das oito células vizinhas.

As regras de evolução são:

\begin{itemize}
    \item uma célula viva com menos de duas vizinhas vivas morre;
    \item uma célula viva com duas ou três vizinhas vivas sobrevive;
    \item uma célula viva com mais de três vizinhas vivas morre;
    \item uma célula morta com exatamente três vizinhas vivas torna-se viva.
\end{itemize}

\section{Entrada e saída}

A aplicação recebe o caminho de um arquivo contendo a quantidade de gerações,
as dimensões da matriz, a quantidade inicial de células vivas e suas
coordenadas. O estado final do mundo é escrito na saída padrão no mesmo
formato definido para a entrada.

As coordenadas são interpretadas como linha e coluna, com índices iniciados em
zero. Células localizadas fora dos limites da matriz são consideradas mortas.
Entradas com dimensões inválidas, coordenadas repetidas ou coordenadas fora da
matriz são rejeitadas.

\section{Implementação sequencial}

A versão sequencial representa a matriz em um vetor contíguo de valores
booleanos. O índice de uma célula é calculado a partir de sua linha, de sua
coluna e da quantidade de colunas da matriz.

Para cada célula, o programa percorre as oito posições vizinhas e contabiliza
somente aquelas que estão dentro dos limites da matriz. Em seguida, aplica as
regras do Jogo da Vida para determinar o estado da célula na próxima geração.

A simulação utiliza duas matrizes: uma para leitura da geração atual e outra
para escrita da próxima geração. Ao final de cada geração, os ponteiros são
trocados. Dessa forma, nenhuma célula utiliza estados que já tenham sido
atualizados na mesma geração.

O kernel que calcula um intervalo de linhas está em
\path{src/simulacao.h} e é compartilhado com a versão Pthreads. A referência
sequencial aplica esse kernel a todas as linhas, evitando que diferenças no
cálculo das células sejam confundidas com ganhos de paralelização.

Depois da última geração, o programa conta as células vivas e percorre a matriz
em ordem de linhas e colunas para escrever suas coordenadas na saída padrão.

\section{Perfil de desempenho sequencial}

O perfil da versão sequencial é obtido com \texttt{gprof}, Cachegrind e Massif.
Para o \texttt{gprof}, a versão sequencial é instrumentada com \texttt{-pg}; o
inlining e a substituição interprocedural de argumentos são desativados para
permitir distinguir as funções no perfil. O Cachegrind analisa referências e
faltas nos caches simulados, e o Massif acompanha o uso de memória. A coleta
detalhada também registra a topologia da CPU, a hierarquia de caches e o código
assembly produzido pelo compilador.

O script \texttt{scripts/coletar\_perfil\_sequencial.sh} coleta Cachegrind,
Massif, topologia e assembly. O alvo
\texttt{profile-sequential-gprof} usa um executável instrumentado separado e
grava o perfil em \path{resultados/perfil_sequencial/gprof.txt}. Assim, as
opções de profiling não alteram os binários usados nos experimentos.

O perfil \texttt{gprof} obtido com \texttt{entradas/classicos/acorn.txt}, uma
matriz de $100\times100$ e $500$ gerações, registrou $5$ milhões de chamadas ao
kernel de evolução. \texttt{contar\_vizinhas} concentrou $56{,}52\%$ do tempo próprio,
enquanto \texttt{obter\_indice} concentrou $39{,}13\%$. Juntas, essas funções
representaram a maior parte do custo sequencial. Esse resultado é coerente com
a natureza de acesso irregular da vizinhança: a implementação percorre oito
células por posição e chama uma função para calcular o índice linear em cada
vizinho.

A oportunidade direta de otimização é specializar o tratamento das células
internas e das bordas, evitando verificações repetidas dos limites em
\texttt{contar\_vizinhas}. A instrução usa uma matriz contígua, o que permite
manter a localidade das linhas horizontais. Uma implementação otimizada deve
preservar o kernel comum ou ser validada novamente contra a referência.

\section{Oportunidades de paralelização}

A principal oportunidade de paralelização está na atualização das células da
matriz. Como cada célula da próxima geração depende apenas do estado imutável da
geração atual, diferentes regiões podem ser calculadas simultaneamente.

A matriz é armazenada por linhas; por isso, a decomposição em faixas horizontais
contíguas favorece a localidade de memória. As escritas são independentes, pois
cada participante atualiza apenas sua faixa na matriz da próxima geração. As
leituras podem alcançar faixas vizinhas, mas utilizam somente a matriz da geração
atual.

\section{Estratégia de paralelização com Pthreads}

A matriz é dividida estaticamente em faixas horizontais. O quociente da divisão
da quantidade de linhas por $P$ define o tamanho base, e as linhas restantes são
distribuídas entre os primeiros participantes. Assim, a diferença entre a maior
e a menor faixa é de no máximo uma linha.

O valor $P$ representa o total de threads participantes do cálculo, incluindo a
thread principal. O programa cria $P-1$ workers persistentes, enquanto a própria
thread principal assume a faixa restante. Duas matrizes booleanas contíguas são
compartilhadas: uma é
lida como geração atual e a outra recebe a próxima geração. Cada participante
escreve somente em sua própria faixa. Uma barreira reutilizável ao final de cada
geração garante que nenhum participante troque os papéis das matrizes enquanto
outro ainda calcula a geração anterior.

A thread principal cria as threads auxiliares, libera todas para o cálculo por
meio de uma variável de condição e processa sua própria faixa. Ao final, aguarda
os $P-1$ workers com \texttt{pthread\_join}. O controle de início impede que uma
criação parcial de threads conduza alguns participantes a uma barreira que
nunca seria completada. Nesse caso, as threads já criadas são liberadas em modo
de cancelamento e os recursos são recolhidos.

A quantidade total de participantes deve satisfazer $1 \leq P \leq L$, em que
$L$ é a quantidade de linhas da matriz. Para $P=1$, a simulação ocorre
inteiramente na thread principal, sem criar uma thread auxiliar. Portanto, essa
configuração mede a infraestrutura da versão Pthreads sem concorrência entre
threads.

As duas versões chamam o mesmo kernel para evoluir uma faixa de linhas. A
diferença experimental fica concentrada na criação das threads auxiliares, na
decomposição do domínio e na sincronização por barreira.

\section{Avaliação de desempenho}

\subsection{Métricas}

A avaliação comparará os tempos médios internos das versões sequencial e
paralela para a mesma entrada e no mesmo ambiente. O \textit{speedup} em relação
à implementação sequencial será calculado por

\begin{equation}
    S(P) = \frac{T_{\mathrm{sequencial}}}{T_{\mathrm{paralelo}}(P)},
\end{equation}

em que $P$ representa o total de participantes. A eficiência correspondente
será calculada por

\begin{equation}
    E(P) = \frac{S(P)}{P}.
\end{equation}

Para separar o ganho do kernel Pthreads do ganho causado pelo aumento de
participantes, também será calculada a escalabilidade relativa a uma thread:

\begin{equation}
    S_{1}(P) = \frac{T_{\mathrm{Pthreads}}(1)}{T_{\mathrm{Pthreads}}(P)},
    \qquad
    E_{1}(P) = \frac{S_{1}(P)}{P}.
\end{equation}

O custo e o sobrecusto paralelo serão obtidos, respectivamente, por

\begin{equation}
    C(P) = P \cdot T_{\mathrm{paralelo}}(P),
    \qquad
    O(P) = C(P) - T_{\mathrm{sequencial}}.
\end{equation}

Essas duas referências não são intercambiáveis. A primeira compara duas
implementações distintas; a segunda mostra somente a evolução da implementação
Pthreads quando o número de participantes aumenta.

\subsection{Conjunto de entradas e configurações}

Os experimentos percorrem 18 entradas clássicas, 35 entradas pseudoaleatórias e
seis cargas de desempenho, com matrizes quadradas e retangulares de diferentes
dimensões, densidades e quantidades de gerações. A série local utiliza cinco
repetições por padrão. No GitHub Actions, os valores padrão são dez repetições
e $P \in \{1,2,4,8\}$, com parâmetros configuráveis antes da execução.

Configurações em que $P>L$ não são executadas, pois cada participante deve
receber ao menos uma linha. As duas versões são compiladas em C11 com otimização
\texttt{-O2} dentro do mesmo ambiente baseado em Ubuntu 24.04. O compilador, o
sistema operacional, a afinidade permitida, a topologia da CPU e a memória
disponível são registrados junto aos dados brutos.

O planejamento contempla campanhas em ambientes com dois e quatro
núcleos físicos disponíveis. Em cada ambiente, a referência sequencial e
todas as configurações Pthreads devem ser medidas na mesma execução. Tempos
absolutos de máquinas diferentes não serão combinados no cálculo de speedup.

\subsection{Protocolo de medição}

A versão sequencial é medida na mesma série e na mesma máquina das configurações
paralelas. Cada rodada percorre todas as entradas antes do início da próxima, e
o ponto inicial e o sentido do percurso variam entre as rodadas. Para cada
entrada, as configurações também são rotacionadas, distribuindo as versões e as
quantidades de threads entre diferentes posições da sequência. Em um ciclo
completo, cada configuração ocupa cada posição uma vez; o ciclo seguinte utiliza
o sentido inverso.

Antes de cada processo medido, um programa auxiliar percorre um buffer quatro
vezes maior que a maior cache de dados detectada, limitado entre 64 MiB e
512 MiB. Esse condicionamento ocorre fora da região temporizada e reduz a
dependência da cache compartilhada deixada pela execução anterior.

O tempo principal, registrado em \texttt{tempo\_simulacao\_segundos}, utiliza
\texttt{clock\_gettime} com \texttt{CLOCK\_MONOTONIC}. A leitura da entrada ocorre
antes dessa região e a escrita do estado final é desativada durante os
experimentos. A região temporizada inclui a alocação da segunda matriz nas duas
versões. Na versão paralela, inclui também a inicialização dos mecanismos de
sincronização, a criação das threads auxiliares, as barreiras, os
\texttt{joins} e a liberação desses recursos.

As métricas adicionais coletadas por \texttt{/usr/bin/time}---tempo de usuário,
tempo de sistema, tempo de parede, CPU, memória, paginação, trocas de contexto e
E/S---abrangem o processo completo. Elas servem como caracterização auxiliar e
não substituem o tempo interno utilizado no cálculo de speedup e eficiência.

O CSV registra a ordem global, a posição da configuração e a execução
precedente. A independência absoluta não é garantida, pois frequência,
temperatura, escalonamento e outras atividades do sistema permanecem
compartilhados. Seus efeitos são reduzidos pelas repetições, pela ordem
balanceada e pelo uso conjunto de média, mediana e dispersão. A cache de arquivos
não é descartada, porque a leitura não participa do tempo da simulação e sua
limpeza alteraria globalmente a máquina.

Além das médias, são calculados desvio-padrão amostral, mínimo, máximo e
mediana. A mediana complementa a média por ser menos sensível a medições
atípicas. Configurações com mais participantes que núcleos físicos disponíveis
são marcadas como \textit{oversubscription} e devem ser interpretadas
separadamente.

\subsection{Validação funcional}

A correção funcional é uma condição anterior à análise de desempenho. A
automação compara as saídas das versões sequencial e Pthreads em casos de borda,
divisões regulares e irregulares das linhas e entradas pseudoaleatórias. Também
verifica a rejeição de entradas e parâmetros inválidos.

Uma validação local reduzida exercita especificamente a participação da thread
principal com $P \in \{1,2,4\}$, quantidades pares e ímpares de gerações e um
limite de tempo para detectar possíveis bloqueios. A mesma validação pode ser
executada sob o Memcheck para procurar acessos inválidos e vazamentos de
memória, mantendo o limite de tempo como proteção contra bloqueios.

\subsection{Ameaças à validade}

O condicionamento da cache reduz, mas não elimina, a influência da execução
anterior. Frequência dinâmica, temperatura, escalonamento do sistema operacional
e atividades externas ainda podem alterar as medições. Em runners
virtualizados, a máquina física e a contenção com outros hóspedes também podem
não ser controláveis.

As repetições, a rotação da ordem e o registro do ambiente reduzem esses riscos,
mas não os anulam. Por isso, diferenças pequenas devem ser analisadas junto com
a dispersão, e execuções com \textit{oversubscription} não devem ser tratadas
como evidência direta da escalabilidade disponível em hardware com mais
processadores.

\subsection{Automação e artefatos}

O script \path{scripts/executar_experimentos_completos.sh} produz o CSV com
todas as execuções, o resumo de desempenho e a descrição do ambiente em
\path{resultados/experimentos_completos/}. O relatório Markdown é gerado por
\path{scripts/gerar_relatorio_experimentos.sh}.

O workflow \texttt{experimentos-completos.yml} compila as duas versões, compara
as saídas em casos de borda e entradas pseudoaleatórias, verifica a rejeição de
entradas inválidas, executa os experimentos, coleta Cachegrind, Massif e
\texttt{gprof}, e publica o relatório e os dados brutos como artefato. Esse fluxo
mantém a referência sequencial e as medições Pthreads no mesmo runner.

\section{Resultados e discussão}

A campanha local foi concluída em 23 minutos e 38 segundos em Ubuntu 24.04,
Linux 6.17, com um Intel Core i5-11300H, quatro núcleos físicos e oito
processadores lógicos. Foram registradas 292 execuções: 59 versões sequenciais e 233 configurações
Pthreads. A cache foi condicionada com um buffer de 64 MiB antes de cada
medição. Como houve apenas uma repetição nesta execução local, os resultados
devem ser interpretados como uma caracterização inicial, não como uma campanha
estatística com múltiplas repetições.

A tabela abaixo sintetiza duas cargas de desempenho representativas.

\begin{table}[htbp]
\centering
\begin{tabular}{c|rrrrr}
\hline
\textbf{Carga} & $P$ & $T_{seq}$ (s) & $T_P$ (s) & $S$ & $E$\\
\hline
$1024^2$, 500 gerações & 1 & 28,278846 & 19,362006 & 1,460533 & 146,05\%\\
$1024^2$, 500 gerações & 2 & 28,278846 & 13,070541 & 2,163556 & 108,18\%\\
$1024^2$, 500 gerações & 4 & 28,278846 & 8,318745 & 3,399413 & 84,99\%\\
$1024^2$, 500 gerações & 8 & 28,278846 & 7,300683 & 3,873452 & 48,42\%\\
\hline
$8192^2$, 50 gerações & 1 & 145,074737 & 104,782795 & 1,384528 & 138,45\%\\
$8192^2$, 50 gerações & 2 & 145,074737 & 63,034194 & 2,301524 & 115,08\%\\
$8192^2$, 50 gerações & 4 & 145,074737 & 43,686934 & 3,320781 & 83,02\%\\
$8192^2$, 50 gerações & 8 & 145,074737 & 31,999475 & 4,533660 & 56,67\%\\
\hline
\end{tabular}
\caption{Tempos e ganhos das cargas quadradas de maior duração.}
\label{tab:desempenho_quadrado}
\end{table}

Nas duas cargas, a execução com dois participantes foi aproximadamente de
$2{,}16$ a $2{,}30$ vezes mais rápida que a sequencial. Com quatro, o ganho
foi próximo de $3{,}32$ a $3{,}40$ vezes. Com oito participantes, o melhor
resultado observado foi de $4{,}53$ vezes na carga $8192^2$, com eficiência de
$56{,}67\%$. Esse comportamento é compatível com quatro núcleos físicos e oito
processadores lógicos. As configurações acima de quatro participantes são
classificadas como oversubscription neste hardware e incluem simultaneamente o
custo adicional de criação das threads, barreiras e trocas de contexto.

Os resultados retangulares também confirmaram a importância da carga. Na
matriz horizontal $2048\times8192$, o tempo caiu de $101{,}600184$ s em uma
thread para $28{,}496014$ s com oito, obtendo speedup de $3{,}565418$ e
eficiência sequencial de $44{,}57\%$. Na matriz vertical $8192\times2048$, a
melhoria foi de $103{,}992752$ s para $21{,}569500$ s, com speedup de
$4{,}821287$ e eficiência de $60{,}27\%$. Não é correto comparar diretamente
matrizes retangulares com a carga quadrada, mas a tendência mostra que o ganho
cresce quando existe trabalho suficiente para amortizar a sincronização.

Em cargas pequenas, como matrizes $5\times5$, $10\times10$ e $20\times20$, o
Pthreads de uma thread pode ser mais rápido ou mais lento que a implementação
sequencial devido às diferenças de infraestrutura. Ao aumentar a quantidade
de participantes, os tempos frequentemente aumentam ou apresentam oscilações,
porque o custo de barreiras e criação supera o trabalho disponível em cada
faixa. Isso explica por que a avaliação deve separar cargas pequenas e cargas
de desempenho.

Não foram realizadas variações controladas de dois ambientes com disponibilidade
diferente de processadores nesta campanha local. Portanto, não há evidência
para confirmar ou rejeitar a escalabilidade além de oito participantes. Como
cada carga foi executada uma vez, os picos numéricos devem ser interpretados
com cautela. A campanha possui repetição configurada na automação e pode ser
executada novamente com cinco ou dez repetições para obter desvios, medianas e
intervalos estatísticos.

\section{Considerações finais}

A implementação sequencial e paralela preserva os resultados em 47 casos de
equivalência funcional e rejeita 15 classes de entradas inválidas. A thread
principal participa da decomposição em todas as configurações testadas.

O gprof identificou 	exttt{contar\_vizinhas} e 	exttt{obter\_indice} como os
principais consumidores do tempo sequencial. A paralelização foi portanto bem
adequada ao domínio, pois a atualização das células é independente dentro de
cada faixa. Os resultados reais mostram ganho em cargas grandes, eficiência
razoável até quatro threads e queda de eficiência em oito threads, coerente com
o custo da sincronização e com oito threads em oversubscription sobre quatro
núcleos físicos.

A implementação atende à interface definida no enunciado: entrada textual,
gerações, dimensões, células iniciais e coordenadas, além da saída no mesmo
formato. Os dados brutos, o resumo, o perfil gprof e o relatório Markdown estão
presentes no diretório 	exttt{resultados}. Valgrind não estava instalado no
environment local sem privilégios administrativos, portanto Cachegrind e Massif
não puderam ser executados nesta máquina.
\end{document}
